
\begin{frame}[t]{The control room receives a recommendation with an evidence trail}

\begin{columns}[T]
\column{.51\textwidth}
\includegraphics[width=\linewidth,height=.56\textheight,keepaspectratio]{\figfile{proposal_control_room.png}}
\scriptsize Concept illustration; AI-generated, not a deployed control room.
\column{.46\textwidth}
\textbf{What the operator needs}\par
Requested action and exact equipment or candidate ID.\par\medskip
Current case, model scope and constraint margins.\par\medskip
Failed alternatives, unresolved assumptions and age of evidence.\par\medskip
Concise public rationale, with each numerical claim tied to a tool result.
\end{columns}
\par\smallskip\small Deployment concept: advisory service beside existing operational applications; authorized procedures govern any later field action.

\slidetransition{How does a review become an explicit acknowledgment?}
\sources{\href{https://arxiv.org/abs/2609.22476}{Governance-aware TSO digital twin (2026)}; illustration provenance in \texttt{figures/Proposal\_Illustration\_Provenance.md}.}
\note{162--164 min. Explain the deployment boundary: notebook, numerical sandbox and advisory review are implemented; operational integration is a concept. Do not suggest the synthetic feeder is a live network. Ask what the operator would need to see before reviewing a change during a busy shift. Spoken bridge: How does a review become an explicit acknowledgment?}
\end{frame}

\begin{frame}[t]{Lab 12: acknowledgment is a state transition, not a generated sentence}

\checkpoint{L12}{Build an operator packet and test the readback gate.}
\cellcode{L12}{9}{10}
\exercise{4}{Try a wrong candidate ID, then the correct readback. Change the requested control after review. Decide whether the old acknowledgment should remain valid.}
\begin{center}\begin{tikzpicture}[every node/.style={font=\small,align=center},b/.style={draw=forest,text width=2.55cm,minimum height=.9cm}]
\node[b] (a) at (0,0) {Proposed};\node[b] (b) at (3.4,0) {Verified};\node[b] (c) at (6.8,0) {Awaiting review};\node[b,fill=mint!25] (d) at (10.2,0) {Acknowledged};
\uncover<2->{\draw[->,thick,forest] (a)--(b);}\uncover<3->{\draw[->,thick,forest] (b)--(c);}\uncover<4->{\draw[->,thick,forest] (c)--(d);}
\end{tikzpicture}\end{center}
\small This is a local advisory exercise. No state in this notebook dispatches equipment.

\slidetransition{Can we reconstruct every stage from one saved trace?}
\sources{Executed classroom experiment: \texttt{Classroom\_Labs.ipynb}; source: \texttt{src/classroom\_labs.py}. Teaching assumptions are explicit. Runtime authority: \href{https://arxiv.org/abs/2609.22476}{Mylonas et al. (2026)}.}
\note{164--168 min. Four minutes hands-on. The operator reviews a specific packet; a language model cannot acknowledge on their behalf. Wrong identity is rejected. If a changed action is not reverified, a prior approval cannot carry over. A completed readback is still not proof the operator understood the evidence; that needs human-factors evaluation. Spoken bridge: Can we reconstruct every stage from one saved trace?}
\end{frame}

\begin{frame}[t]{One trace should answer what happened at every stage}

\checkpoint{L10}{Compare real saved stage records and public responses.}
\cellcode{L10}{1}{3}\footnotesize
\begin{tabular}{@{}>{\raggedright\arraybackslash}p{.22\linewidth}>{\raggedright\arraybackslash}p{.69\linewidth}@{}}
\toprule Inspect & What the answer should expose \\
\midrule
Input & Case, measurement assumptions, allowed operations and model scope. \\
Method & Which optimizer, policy, simulator or language call actually ran. \\
Output & Candidate ID, returned values, constraint failures and public reason. \\
Cost & Measured tool/call time, including failed attempts. \\
Disposition & Retained, rejected, refined or escalated; the next required step. \\
\bottomrule
\end{tabular}\par
\par\medskip
\textbf{Trace labels matter:} classroom fresh solve, archived actual LLM response and optional live call are different evidence modes.

\slidetransition{Which missing evidence changes the operator decision?}
\sources{Executed classroom experiment: \texttt{Classroom\_Labs.ipynb}; source: \texttt{src/classroom\_labs.py}. Teaching assumptions are explicit. Actual recorded five-stage reference: \texttt{Proposal\_Complete\_Workflow.ipynb}, P1--P8; \texttt{results/proposal\_complete\_workflow.json}.}
\note{168--170 min. L10 is archived five-stage evidence, clearly labeled. The fresh classroom methods are related but are not retroactively given those old language responses. Ask students to find a stage that uses a deterministic method and explain why it still belongs in the agentic workflow. Each role has a responsibility, not necessarily its own model process. Spoken bridge: Which missing evidence changes the operator decision?}
\end{frame}

\begin{frame}[t]{Capstone: would you approve, hold or request another study?}

\begin{tabular}{@{}>{\raggedright\arraybackslash}p{.13\linewidth}>{\raggedright\arraybackslash}p{.57\linewidth}>{\raggedright\arraybackslash}p{.22\linewidth}@{}}
\toprule Card & Evidence presented & Your decision \\
\midrule
A & MV PASS; full-feeder AC FAIL; Llama says ``ready''. & Review / hold / study? \\
B & Restoration MILP PASS; equipment isolation has no supporting record. & Review / hold / study? \\
C & No residual alarm; protected measurement set is unavailable. & Review / hold / study? \\
\bottomrule
\end{tabular}\par
\par\medskip
\textbf{Task:} identify the missing fact and the appropriate tool or human step. Do not answer with a model's confidence score.
\uncover<2->{\par\medskip\textbf{Debrief:} A requires broader-model repair; B requires authorized isolation evidence; C requires justified independent information or escalation. None supports automatic dispatch.}

\slidetransition{Which architecture fits each kind of operator task?}
\sources{Executed classroom experiment: \texttt{Classroom\_Labs.ipynb}; source: \texttt{src/classroom\_labs.py}. Teaching assumptions are explicit. Schweppe \& Wildes, \href{https://doi.org/10.1109/TPAS.1970.292678}{state estimation (1970)}; Liu, Ning \& Reiter, \href{https://doi.org/10.1145/1653662.1653666}{false-data injection, CCS 2009}.}
\note{170--172 min. Ask one row per group. These cards are intentionally distinct: model scope, physical permission and information trust. The correct answer includes a concrete next action rather than simply no. Do not imply that holding a prior setting is always safe after a topology or measurement change. Spoken bridge: Which architecture fits each kind of operator task?}
\end{frame}

\begin{frame}[t]{Match the architecture to the operator task}

\begin{tabular}{@{}>{\raggedright\arraybackslash}p{.30\linewidth}>{\raggedright\arraybackslash}p{.31\linewidth}>{\raggedright\arraybackslash}p{.31\linewidth}@{}}
\toprule Task condition & Useful architecture & Evidence required \\
\midrule
Stable repeated calculation & Fixed workflow + numerical checks & Correct model and tested contract. \\
Varied language requests & LLM tool selection + typed state & Intent accuracy and supported claims. \\
Several independent studies & Bounded specialist roles & Handoff quality and real parallel benefit. \\
Competing refinements & Vote or learned selector & Held-out gain at equal trial budget. \\
High-consequence action & Human review + enforced authorization & Correct action, scope, provenance and permissions. \\
\bottomrule
\end{tabular}\par
\par\medskip
\textbf{Design principle:} add complexity only when it solves a measured failure of the simpler workflow.

\slidetransition{Can you explain one complete case using the five stages?}
\sources{\href{https://arxiv.org/abs/2512.08296v3}{Kim et al., scaling agent systems (2026)}; \href{https://arxiv.org/abs/2609.22476}{governance-aware TSO study (2026)}.}
\note{172--174 min. This table is a synthesis, not a ranking of packages. Ask learners to choose the simplest design for their own use case. A deterministic workflow can include the same engineering verifier and operator controls; those features are not exclusive benefits of an LLM. Spoken bridge: Can you explain one complete case using the five stages?}
\end{frame}

